\section{Introduction}
\label{sec:introduction}

Electronic medical records (EMRs) enable health information to follow a patient across services, but integration also enlarges the set of parties and systems that may process sensitive clinical data. Indonesian regulation requires healthcare facilities to implement integrated EMRs, while health information is classified as specific personal data that must be protected from unauthorized processing \cite{permenkes24_2022,uu27-2022-pdp}. Consequently, an integrated EMR system must support data availability for care without turning a valid clinical relationship into unrestricted access.

Patient-centric access control gives patients authority over who may use their records \cite{peng2023patientcentric}. Direct care, however, is collaborative: medical-records personnel, nurses, physicians, laboratory personnel, and pharmacists participate at different stages and require different subsets of the record \cite{who2010_interprofessional_collaborative_practice}. The Indonesian EMR metadata guideline similarly separates outpatient, inpatient, laboratory, and pharmacy datasets and distinguishes administrative from clinical documentation \cite{kemenkes1423-2022}. These workflows require access decisions over the actor, operation, episode, data category, and time---not merely a distinction between administrative and medical personnel.

DecMed is a patient-centric decentralized EMR prototype that combines IOTA for on-chain state, IPFS for encrypted off-chain records, and proxy re-encryption (PRE) for controlled data disclosure \cite{sadewa2025decmed}. Its previous capability token was a JSON Web Token (JWT) containing an IOTA address, a broad role, and a read/update purpose. The protected clinical record was also represented as a relatively large encrypted object. This combination created three connected limitations. First, access issuance still depended on the patient for each recipient. Second, the policy was coarse-grained because decisions were driven mainly by role and purpose. Third, the storage unit did not provide segments on which a finer policy could be enforced. The bearer-token workflow also needed evidence that the requester controlled the wallet associated with the active actor, together with revocation and audit semantics for derived grants.

This paper extends DecMed with a Policy-Based Access Control (PBAC) model, segmented EMR objects, and Macaroon capability tokens. PBAC defines the relationships among actors, EMR objects, operations, and conditions. Macaroons were selected for their append-only caveats and attenuation property, rather than on a claim that they are universally more secure than JWT, PASETO, or Biscuit. A patient still authorizes the initial access. A legitimate holder can subsequently request a narrower child token from the PRE Server, which validates the parent capability before producing the attenuation result. Macaroons conceptually allow attenuation without involving the original issuer; in this implementation, child-token formation is mediated by the PRE Server so that parent token, delegator identity, requested scope, and delegation chain are checked before issuance. Delegation therefore does not require the patient's direct involvement at every handover among healthcare actors, but it still requires communication with the PRE Server.

The contributions are:
\begin{enumerate}
	\item a PBAC-based authorization model for DecMed, represented through Macaroon caveats for read/update scope, episode, hospital, expiration, delegation depth, and parent--child attenuation;
	\item an EMR segmentation model that aligns authorization with combinations of \texttt{dataset\_category} and \texttt{function\_category};
	\item PRE-mediated token attenuation for delegation among healthcare actors, integrated with active-wallet proof, institutional role checks, Redis-backed revocation, and IOTA audit records; and
	\item a prototype evaluation covering collaborative outpatient and inpatient workflows and negative authorization cases.
\end{enumerate}

Section~\ref{sec:background} summarizes the technical context and related work. Section~\ref{sec:approach} develops the problem analysis and proposed approach. Sections~\ref{sec:architecture} and \ref{sec:implementation} describe the system and its implementation. Sections~\ref{sec:evaluation} and \ref{sec:discussion} present the evidence and limitations, followed by the conclusion.
